\begin{table}[th]
  \caption{Comparison between different proposed interfaces with HuBERT Base as the upstream model, evaluated on the ML-SUPERB Monolingual track. ``GroupWS'' indicates Grouped Weighted Sums and ``Hierarchical Conv.'' stands for Hierarchical Convolution over layers. The metric is WER/CER.}
  \label{tab:preliminary}
  \centering
  \begin{tabular}{ l c c c }
    \toprule
    Interface & Params & Mono-10min & Mono-1hr \\
    \midrule
    Weighted Sum & 13 & 42.85           & 35.15 \\
    \midrule
    GroupWS(\#Groups=2) & 1.2M & 41.84     & 34.47 \\
    GroupWS(\#Groups=3) & 1.8M & 43.08     & 33.96 \\
    GroupWS(\#Groups=4) & 2.4M & 42.52     & 33.99 \\
    Concat $+$ Proj & 7.7M & 43.20      & 34.26 \\
    % Transform $+$ WS & 7.7M & 42.42     & 34.97 \\
    PCA $+$ Concat & 0 & 45.16          & 36.58 \\
    Hierarchical Conv. & 4.4M & \textbf{41.51}   & \textbf{33.88} \\
    CLS Pooling & 5.5M & 48.36          & 33.92 \\

    \bottomrule
  \end{tabular}
  \vspace{-15pt}
\end{table}


\begin{table*}[th]
  \caption{Performance comparison of proposed interfaces using different upstream models, evaluated on both ML-SUPERB and SUPERB. The metric for each task: 
  ``Mono-1hr'' and ``Multi-1hr'': WER/CER,
  ``LID-1hr'',``ER'',``IC'', ``SV'': accuracy,
  ``PR'': PER.
  % \ian{The table seems to be oversized. Perhaps we should remove Mono-10min?}
  }
  \label{tab:full_results}
  \centering
  \begin{tabular}{ l l c c c c c c c }
    \toprule
    \multirow{ 2}{*}{Upstream}  & \multirow{ 2}{*}{Interface} & \multicolumn{3}{c}{ML-SUPERB} & \multicolumn{4}{c}{SUPERB}   \\
    % \cmidrule{4-6} 
    &  & Mono-1hr~($\downarrow$) & Multi-1hr~($\downarrow$)  & LID-1hr~($\uparrow$)
    & ER~($\uparrow$) & IC~($\uparrow$) & SV~($\downarrow$) & PR~($\downarrow$) \\
    \midrule
    \multirowcell{ 3}{HuBERT\\Base} 
    &  Weighted Sum  &	35.1 & 24.4 & \textbf{84.6} & 64.92 & 98.34 & \textbf{5.11} & 5.41 \\
    &  Hierarchical Conv.  &	\textbf{33.9} & \textbf{24.0} & 81.7 & \textbf{68.49} & \textbf{99.45} & 5.62 & 3.07\\
    &  CLS Pooling &	\textbf{33.9} & 24.3 & 1.2 & 62.73 & \textbf{99.45} & 49.91 & \textbf{2.93} \\
    \midrule

    \multirowcell{ 3}{HuBERT\\Large} 
    &  Weighted Sum         &	32.3 & 22.3 & 64.4 & 67.62 & 98.76 & \textbf{5.98} & 3.53 \\
    &  Hierarchical Conv.    &	\textbf{30.0} & \textbf{21.4} & \textbf{89.2} & \textbf{72.44} & \textbf{99.53} & 6.03 & 1.76 \\
    &  CLS Pooling           &	31.1 & 21.5 & 87.0 & 68.72 & \textbf{99.53} & 6.34 & \textbf{1.64} \\
    \midrule



    \multirowcell{ 3}{WavLM\\Base} 
    &  Weighted Sum  &	34.2 & 24.3 & \textbf{84.1} & 65.94 & 98.63 & \textbf{4.69} & 4.84 \\
    &  Hierarchical Conv.  &	\textbf{32.4} & \textbf{23.6} & 67.9 & \textbf{68.57} & \textbf{99.53} & 5.48 & 3.06\\
    &  CLS Pooling  &	33.8 & 23.8 & 13.2 & 63.60 & 99.50 & 6.03 & \textbf{2.92}\\
    \midrule

    \multirowcell{ 3}{WavLM\\Large} 
    &  Weighted Sum  &	30.1 & 20.8 & 71.7 & 70.62 & 99.31 & \textbf{3.77} & 3.06\\
    &  Hierarchical Conv.  &	\textbf{28.0} & 19.4 & \textbf{90.6} & \textbf{74.95} & \textbf{99.71} & 5.20 & \textbf{1.72}\\
    &  CLS Pooling &	29.7 & \textbf{18.8} & 88.5 & 70.45 & 99.63 & 4.27 & 1.73 \\
    \midrule

    \multirow{ 3}{*}{XLSR-53} 
    &  Weighted Sum  &	35.1 & 20.2 & 76.6 & 66.34 & 95.62 & 6.45 & 4.50\\
    &  Hierarchical Conv.  &	\textbf{30.6} & \textbf{19.8} & \textbf{80.4} & \textbf{72.01} & \textbf{99.55} & \textbf{5.63} & \textbf{2.69}\\
    &  CLS Pooling &	31.0 & 80.7 & 4.5 & 65.10 & 99.47 & 12.95 & 5.27\\   

    \bottomrule
  \end{tabular}
  \vspace{-10pt}
\end{table*}



\subsection{Experimental Setup}

To get a holistic performance comparison of different interface designs, we choose 5 SSL models and 2 speech SSL benchmarks for evaluation.
Specifically, we pick HubERT Base and Large~\cite{hsu2021hubert}, WavLM Base and Large~\cite{Chen2021WavLMLS}, and XLSR-53~\cite{Conneau2020xlsr} for our upstream models.
For downstream tasks, we test all proposed interfaces on ML-SUPERB~\cite{shi23g_mlsuperb} because it is currently less saturated than many SUPERB tasks.
ML-SUPERB contains both 10min and 1hr training sets for each language.
In the ML-SUPERB Monolingual track, the overall reported score is the average CER across 13 languages (the model for each language is trained individually).
In the multilingual track, there are two tasks: \textbf{ASR} and Language Identification~(\textbf{LID}).
In these tasks, a model must be trained on 143 languages simultaneously.

In addition to multilingual ASR tasks, we also test the generalization of our interface designs on other common speech processing tasks.
Hence, we select 4 tasks from  SUPERB: Intent Classification~(\textbf{IC}), Phoneme Recognition~(\textbf{PR}), Emotion Recognition~(\textbf{PR}), and Speaker Verification~(\textbf{SV}).

To reduce the computational cost of running a large number of potential experiments, we first test all proposed interfaces on ML-SUPERB Monolingual Track with HuBERT Base.
Then we pick the top 2 most promising interface designs and run them on all 5 upstream models and both ML-SUPERB and SUPERB.
For both SUPERB and ML-SUPERB, we follow the default training configurations as far as possible.
However, for the Multilingual track on ML-SUPERB, we find that the default downstream model
% ~(2 layer of Transformer layer) 
is not optimal, and performance can be easily increased by simply adding more layers.
In order to disentangle the effect of a larger downstream model from the interface itself, we do a hyperparameter search on the best number of layers in the downstream model so that it can achieve the lowest CER in terms of the same frozen upstream model.
From our empirical experiment, we found that 65 is the optimal layer number for Multilingual ASR~(about 56M params).

 
\subsection{Interface Comparison on ML-SUPERB Monolingual ASR with HuBERT Base}

The results for our pilot experiment using HuBERT Base on the ML-SUPERB Monolingual task are shown in Table~\ref{tab:preliminary}. Firstly, we notice that there is a significant performance difference between all proposed interfaces relative to the baseline weighted sum method. While several of the proposed interfaces improve over the baseline weighted sum, we see that the Hierarchical Convolution achieves the overall best interface performance on both the 10min and 1hr datasets.
The CLS Pooling interface ranks second on the 1hr task, but performs poorly on 10min, indicating that this method likely requires more finetuning data to reach its potential.
In light of their promising results on the 1hr ML-SUPERB monolingual task, we choose Hierarchical Convolution and CLS Pooling for further experimentation.

\subsection{Interface Comparison on ML-SUPERB and SUPERB}
The results for the full suite of tasks across ML-SUPERB and SUPERB using all upstream models are shown in Table~\ref{tab:full_results}. Again, we observe that the Hierarchical Convolution interface tends to consistently outperform the baseline weighted sum as well as the CLS Pooling interface.
The performance differences are often substantial: for example, simply replacing the weighted sum interface with the Hierarchical Convolution interface reduces HuBERT Base's PER on the SUPERB phone recognition task from 5.41 to 2.93, which is lower than the 3.53 PER achieved by HuBERT Large with the weighted sum interface.
We even find that the CLS Pooling interface with the HuBERT Large model establishes a new SotA on the SUPERB leaderboard for phone recognition at 1.64 PER, beating the WavLM model's 3.06 PER when using the weighted sum.

Generally, the Hierarchical Convolution interface
% Among all the tasks, the two proposed interfaces 
outperform the baseline weighted sum on all tasks except for Speaker Verification with HuBERT and WavLM models, although XLSR-53 with Hierarchical Convolution outperforms the weighted sum.
This might imply that speaker-related information is encoded similarly across different layers in HuBERT and WavLM, causing our proposed interfaces to overfit.

% Interestingly, when using the Hierarchical Convolution we also observe a larger improvement over the basic weighted sum in Large models compared to Base models for HuBERT and WavLM, which supports our hypothesis that the basic weighted sum averages out more information as the number of layers increases.
% This phenomenon can be observed in both ML-SUPERB and SUPERB.
Interestingly, in both benchmarks, we also observe a larger improvement over the basic weighted sum in Large models compared to Base models for HuBERT and WavLM when Hierarchical Convolution is applied, which supports our hypothesis that the basic weighted sum averages out more information as the number of layers increases.
In particular, on the LID task, we see that with Base upstream models the Hierarchical Convolution interface underperforms the basic weighted sum, but when scaling to the Large models (HuBERT Large, WavLM Large, and XLSR-53) the Hierarchical Convolution instead outperforms the basic weighted sum.


\subsection{Interface vs. Larger Downstream Models}

\begin{table}[th]
\scriptsize
\vspace{-5pt}
  \caption{Ablation experiments on trainable parameters. ``WS w/ Large DS'' indicates Weighted Sum with Large Downstream model.}
  \label{tab:ablation}
  \centering
  \begin{tabular}{ l l c c }
    \toprule
    \multirow{ 2}{*}{Upstream}  & \multirow{ 2}{*}{Interface} & \multicolumn{2}{c}{ML-SUPERB~(WER/CER)}    \\
    % \cmidrule{4-6} 
    &  & Mono-1hr~($\downarrow$) & Multi-1hr~($\downarrow$) \\
    \midrule
    \multirowcell{ 2}{HuBERT\\Base} 
    % \multirow{ 2}{*}{Upstream}
    % &  Weighted Sum &           35.1 & 24.4  \\
    &  Hierarchical Conv. &     \textbf{33.9} & \textbf{24.0} \\
    &  WS w/ Large DS &            35.2 & 24.8  \\
    \midrule

    \multirowcell{ 2}{HuBERT\\Large} 
    % &  Weighted Sum         &	32.3 & 22.3  \\
    &  Hierarchical Conv.   &	\textbf{30.0} & \textbf{21.4}  \\
    &  WS w/ Large DS          &	32.9 & 22.1  \\
    \midrule


    \multirowcell{ 2}{WavLM\\Base} 
    % &  Weighted Sum &           34.2 & 24.3  \\
    &  Hierarchical Conv. &     \textbf{32.4} & \textbf{23.6} \\
    &  WS w/ Large DS &            34.4 & 24.3 \\
    \midrule

    \multirowcell{ 2}{WavLM\\Large} 
    % &  Weighted Sum &           30.1 & 20.8 \\
    &  Hierarchical Conv. &     \textbf{28.0} & \textbf{19.4} \\
    &  WS w/ Large DS &            30.5 & 20.8  \\
    \midrule

    \multirow{ 2}{*}{XLSR-53} 
    % &  Weighted Sum &           35.1 & 20.2 \\
    &  Hierarchical Conv. &     \textbf{30.6} & \textbf{19.8} \\
    &  WS w/ Large DS &            35.4 & 20.0 \\   

    \bottomrule
  \end{tabular}
  \vspace{-5pt}
\end{table}




Although there are substantial improvements using the Hierarchical Convolution interface, it is not immediately clear whether the performance gain comes from the interface design or simply the additional trainable parameters.
To this end, we increase the size of the downstream model for Monolingual and Multilingual ASR on ML-SUPERB by an amount roughly the same as the number of parameters in the Hierarchical Convolution interface (denoted as ``WS+Large DS'').

As shown in Table~\ref{tab:ablation}, under all settings, weighted sum with a large downstream model does not lead to better performance than the Hierarchical Convolution interface combined with a smaller downstream model.
% In fact, when using the large downstream model we begin to see evidence of overfitting.
This implies the role of the interface is not simply to add more learnable parameters to the downstream model, but to actually better aggregate information across the layers of the upstream model.
Our findings not only resonate with~\cite{zaiem23b_interspeech} that larger downstream models increase the performance but we also show that a better interface design can get further improvements.

\subsection{Performance under Fine-tuning setting}
We motivated the introduction of the interface module under the assumption that the upstream model would typically remain frozen during the entire downstream task fine-tuning.
However, we wish to empirically verify whether different interfaces can still offer improvements with end-to-end fine-tuning.
Practically, we follow the same pipeline as in our previous experiments except the upstream model is trainable.
% but now allow the upstream model weights to be trainable.
We evaluate on ML-SUPERB Monolingual track 1h using HuBERT Base.

\begin{table}[th]
\scriptsize
\vspace{-5pt}
  \caption{Comparison of Weighted Sum and Hierarchical under end-to-end fine-tuning. The metric is WER/CER.}
  \label{tab:finetune}
  \centering
  \begin{tabular}{ l c c }
    \toprule
     \multirow{ 2}{*}{Interface} & \multicolumn{2}{c}{ML-SUPERB Mono-1hr~($\downarrow$)}    \\
    % \cmidrule{4-6} 
      & Freeze Upstream & Fine-tuned \\
    \midrule
    
    % \multirow{ 2}{*}{Upstream}
    % &  Weighted Sum &           35.1 & 24.4  \\
    Weighted Sum &     35.1 & 31.5 \\
     
     Hierarchical Conv. &  \textbf{33.9} & \textbf{31.1} \\
      WS w/ Large DS &       35.2 & 31.6 \\

    \bottomrule
  \end{tabular}
  \vspace{-5pt}
\end{table}





As shown in Table~\ref{tab:finetune}, end-to-end fine-tuning shortens the gap between different interface designs.
However, the Hierarchical Convolution interface still outperforms the weighted sum baseline.
% These results indicate that considering the interface design is beneficial no matter whether the upstream model is frozen or not during fine-tuning.
These results indicate that the interface design is beneficial regardless of frozen or trainable upstream model.